\begin{textsample}{POS Dim 1 – 2016 – Score 20.00 – 2016/t001086.txt}  \label{ex:f1_pos_053}
In 1956, a documentary by Jacques Cousteau won both \textbf{the} Palme d ' Or and an Oscar award.
This film was called,"Le Monde Du Silence,"or,"\textbf{The} Silent World."\textbf{The} premise of \textbf{the} title was that \textbf{the} \textbf{underwater} world was a quiet world.
We now know, 60 years later, that \textbf{the} \textbf{underwater} world is anything but silent.
Although \textbf{the} sounds are inaudible above water depending on where you are and \textbf{the} time of year, \textbf{the} \textbf{underwater} soundscape can be as noisy as any jungle or rainforest.
Invertebrates like snapping \textbf{shrimp}, \textbf{fish} and marine mammals all use sound.
They use sound to study their \textbf{habitat}, to keep in communication with each other, to navigate, to \textbf{detect} predators and prey.
They also use sound by listening to know something about their environment.
Take, for an example, \textbf{the} Arctic.
It ' s considered a vast, inhospitable place, sometimes described as a desert, because it is so cold and so remote and ice- covered for much of \textbf{the} year.
And despite this, there is no place on Earth that I would rather be than \textbf{the} Arctic, especially as days lengthen and \textbf{spring} comes.
To me, \textbf{the} Arctic really embodies this disconnect between what we see on \textbf{the} \textbf{surface} and what ' s going on \textbf{underwater}.
You can look out across \textbf{the} ice — all white and \textbf{blue} and cold — and see nothing.
But if you could hear \textbf{underwater}, \textbf{the} sounds you would hear would at first amaze and then delight you.
And while your eyes are seeing nothing for kilometers but ice, your ears are telling you that out there are bowhead and beluga \textbf{whales}, walrus and bearded \textbf{seals}. \textbf{The} ice, too, makes sounds.
It screeches and cracks and pops and groans, as it collides and rubs when temperature or currents or winds change.
And under 100 percent \textbf{sea} ice in \textbf{the} dead of winter, bowhead \textbf{whales} are singing.
And you would never expect that, because we humans, we tend to be very visual animals.
For most of us, but not all, our sense of sight is how we navigate our world.
For marine mammals that live \textbf{underwater}, where chemical cues and light transmit poorly, sound is \textbf{the} sense by which they see.
And sound transmits very well \textbf{underwater}, much better than it does in air, so signals can be heard over great distances.
In \textbf{the} Arctic, this is especially important, because not only do Arctic marine mammals have to hear each other, but they also have to listen for cues in \textbf{the} environment that might indicate heavy ice ahead or open water.
Remember, although they spend most of their lives \textbf{underwater}, they are mammals, and so they have to \textbf{surface} to breathe.
So they might listen for thin ice or no ice, or listen for echoes off nearby ice.
Arctic marine mammals live in a rich and varied \textbf{underwater} soundscape.
In \textbf{the} \textbf{spring}, it can be a cacophony of sound.
But when \textbf{the} ice is frozen solid, and there are no big temperature shifts or current changes, \textbf{the} \textbf{underwater} Arctic has some of \textbf{the} lowest ambient noise levels of \textbf{the} world ' s oceans.
But this is changing.
This is primarily due to a decrease in seasonal \textbf{sea} ice, which is a direct result of human \textbf{greenhouse} gas \textbf{emissions}.
We are, in effect, with climate change, conducting a completely uncontrolled experiment with our \textbf{planet}.
Over \textbf{the} past 30 years, areas of \textbf{the} Arctic have seen decreases in seasonal \textbf{sea} ice from anywhere from six weeks to four months.
This decrease in \textbf{sea} ice is sometimes referred to as an increase in \textbf{the} open water season.
That is \textbf{the} time of year when \textbf{the} Arctic is navigable to vessels.
And not only is \textbf{the} extent of ice changing, but \textbf{the} age and \textbf{the} \textbf{width} of ice is, too.
Now, you may well have heard that a decrease in seasonal \textbf{sea} ice is causing a loss of \textbf{habitat} for animals that rely on \textbf{sea} ice, such as ice \textbf{seals}, or walrus, or polar bears.
Decreasing \textbf{sea} ice is also causing increased erosion along coastal villages, and changing prey availability for marine birds and mammals.
Climate change and decreases in \textbf{sea} ice are also altering \textbf{the} \textbf{underwater} soundscape of \textbf{the} Arctic.
What do I mean by soundscape?
Those of us who eavesdrop on \textbf{the} oceans for a living use instruments called hydrophones, which are \textbf{underwater} microphones, and we record ambient noise — \textbf{the} noise all around us.
And \textbf{the} soundscape describes \textbf{the} different contributors to this noise field.
What we are hearing on our hydrophones are \textbf{the} very real sounds of climate change.
We are hearing these changes from three fronts : from \textbf{the} air, from \textbf{the} water and from land.
First : air.
Wind on water creates waves.
These waves make bubbles ; \textbf{the} bubbles break, and when they do, they make noise.
And this noise is like a hiss or a static in \textbf{the} background.
In \textbf{the} Arctic, when it ' s ice-covered, most of \textbf{the} noise from wind doesn ' t make it into \textbf{the} water column, because \textbf{the} ice acts as a buffer between \textbf{the} atmosphere and \textbf{the} water.
This is one of \textbf{the} reasons that \textbf{the} Arctic can have very low ambient noise levels.
But with decreases in seasonal \textbf{sea} ice, not only is \textbf{the} Arctic now open to this wave noise, but \textbf{the} number of storms and \textbf{the} intensity of storms in \textbf{the} Arctic has been increasing.
All of this is raising noise levels in a previously quiet ocean.
Second : water.
With less seasonal \textbf{sea} ice, subarctic species are moving north, and taking advantage of \textbf{the} new \textbf{habitat} that is created by more open water.
Now, Arctic \textbf{whales}, like this bowhead, they have no dorsal \textbf{fin}, because they have \textbf{evolved} to live and swim in ice-covered waters, and having something sticking off of your back is not very conducive to migrating through ice, and may, in fact, be excluding animals from \textbf{the} ice.
But now, everywhere we ' ve listened, we ' re hearing \textbf{the} sounds of \textbf{fin} \textbf{whales} and humpback \textbf{whales} and killer \textbf{whales}, further and further north, and later and later in \textbf{the} season.
We are hearing, in essence, an invasion of \textbf{the} Arctic by subarctic species.
And we don ' t know what this means.
Will there be competition for food between Arctic and subarctic animals?
Might these subarctic species introduce diseases or parasites into \textbf{the} Arctic?
And what are \textbf{the} new sounds that they are producing doing to \textbf{the} soundscape \textbf{underwater}?
And third : land.
And by land...
I mean people.
More open water means increased human use of \textbf{the} Arctic.
Just this past summer, a massive cruise ship made its way through \textbf{the} Northwest Passage — \textbf{the} once-mythical route between Europe and \textbf{the} Pacific.
Decreases in \textbf{sea} ice have allowed humans to occupy \textbf{the} Arctic more often.
It has allowed increases in \textbf{oil} and gas exploration and extraction, \textbf{the} potential for commercial \textbf{shipping}, as well as increased tourism.
And we now know that ship noise increases levels of stress \textbf{hormones} in \textbf{whales} and can disrupt feeding behavior.
Air guns, which produce loud, low-frequency"whoomps"every 10 to 20 seconds, changed \textbf{the} swimming and vocal behavior of \textbf{whales}.
And all of these sound sources are decreasing \textbf{the} acoustic space over which Arctic marine mammals can communicate.
Now, Arctic marine mammals are used to very high levels of noise at certain times of \textbf{the} year.
But this is primarily from other animals or from \textbf{sea} ice, and these are \textbf{the} sounds with which they ' ve \textbf{evolved}, and these are sounds that are vital to their very survival.
These new sounds are loud and they ' re alien.
They might impact \textbf{the} environment in ways that we think we understand, but also in ways that we don ' t.
Remember, sound is \textbf{the} most important sense for these animals.
And not only is \textbf{the} physical \textbf{habitat} of \textbf{the} Arctic changing rapidly, but \textbf{the} acoustic \textbf{habitat} is, too.
It ' s as if we ' ve plucked these animals up from \textbf{the} quiet countryside and dropped them into a big city in \textbf{the} middle of rush hour.
And they can ' t escape it.
So what can we do now?
We can ' t decrease wind speeds or keep subarctic animals from migrating north, but we can work on local solutions to reducing human- caused \textbf{underwater} noise.
One of these solutions is to slow down ships that traverse \textbf{the} Arctic, because a slower ship is a quieter ship.
We can restrict access in seasons and regions that are important for \textbf{mating} or feeding or migrating.
We can get smarter about quieting ships and find better ways to explore \textbf{the} ocean bottom.
And \textbf{the} good news is, there are people working on this right now.
But ultimately, we humans have to do \textbf{the} hard work of reversing or at \textbf{the} very least decelerating human-caused atmospheric changes.
So, let ' s return to this idea of a silent world \textbf{underwater}.
It ' s entirely possible that many of \textbf{the} \textbf{whales} swimming in \textbf{the} Arctic today, especially long- lived species like \textbf{the} bowhead \textbf{whale} that \textbf{the} Inuits say can live two human lives — it ' s possible that these \textbf{whales} were alive in 1956, when Jacques Cousteau made his film.
And in retrospect, considering all \textbf{the} noise we are creating in \textbf{the} oceans today, perhaps it really was"\textbf{The} Silent World."Thank you.

% matched lemmas: blue, detect, emission, evolve, fin, fish, greenhouse, habitat, hormone, mate, oil, planet, sea, seal, shipping, shrimp, spring, surface, the, underwater, whale, width
\end{textsample}
